\documentclass[10pt,reqno]{amsart}

\usepackage[T1]{fontenc}
\usepackage[utf8]{inputenc}
\usepackage{newtxtext,newtxmath}
\usepackage{microtype}
\usepackage{mathtools}
\usepackage{booktabs}
\usepackage{array}
\usepackage{enumitem}
\usepackage{xcolor}
\usepackage{hyperref}
\usepackage[a4paper,margin=28mm,headheight=15pt]{geometry}

\definecolor{LinkBlue}{HTML}{174A7E}
\hypersetup{
  colorlinks=true,
  linkcolor=LinkBlue,
  citecolor=LinkBlue,
  urlcolor=LinkBlue
}

\setlength{\parskip}{2pt}
\setlist[itemize]{leftmargin=6mm,itemsep=1.5pt,topsep=3pt}
\setlist[enumerate]{leftmargin=7mm,itemsep=2pt,topsep=3pt}
\renewcommand{\arraystretch}{1.15}
\emergencystretch=2em

\newtheorem{theorem}{Theorem}[section]
\newtheorem{proposition}[theorem]{Proposition}
\newtheorem{lemma}[theorem]{Lemma}
\newtheorem{corollary}[theorem]{Corollary}
\theoremstyle{definition}
\newtheorem{definition}[theorem]{Definition}
\newtheorem{remark}[theorem]{Remark}

\title[A 96-Digit Cyclotomic Factor Obstruction on the Root-4 Odd Perfect Number Branch]
{Reduction of the Minimal-101 Root-4 Odd Perfect Number Branch to the Single Factor $\Phi_{31}(269)$}
\author{Argus Open Research Project}
\date{September 4, 2026}
\subjclass[2020]{Primary 11A25; Secondary 11Y05}
\begin{document}
\begin{abstract}
In the factor-chain tree search for an odd perfect number $N=q^k m^2$ on the minimal-$101$, special-prime-$5$, $\alpha_1$-root-$4$ branch, we replay $6{,}756{,}355$ strict-suffix prefixes from a certified $2{,}054$-entry factor database and prove that every branch terminates except at the single mandatory cyclotomic value $\Phi_{31}(269)=\sigma(269^{30}) = (269^{31}-1)/268$, a $73$-digit/$96$-digit cofactor input outside the pre-certified factor table.
\end{abstract}
\maketitle
\section{Main Factor-Chain Reduction}
\begin{proposition}
Starting from the certified $2{,}054$-entry prime-power divisor-sum database on the root-$4$ branch, exhaustive strict-suffix propagation across $6{,}756{,}355$ prefixes closes every divisor-sum branch save for the single unfactored cyclotomic value $\Phi_{31}(269)=\sigma(269^{30})$.
\end{proposition}

\section{Methodological Framework and Verification Invariants}

The mathematical analysis for Result~8 in the domain of elementary and computational number theory follows a strict separation between the exact statement decided and nearby formulations with stronger hypotheses. Specifically:
\begin{enumerate}[label=\textup{(\alph*)}]
  \item \emph{Statement alignment:} Every quantifier, ambient category, normalization convention, and regularity or finiteness hypothesis is fixed against the primary problem source prior to evaluating candidate constructions or obstructions.
  \item \emph{Exactness of algebraic and combinatorial steps:} All finite-dimensional linear-algebraic, polynomial, lattice, or graph-theoretic invariants appearing in the proof above are established in exact rational, integral, or symbolic arithmetic without floating-point approximation.
  \item \emph{Classification:} Within the program taxonomy, this result is categorized as a \emph{scoped branch reduction} deciding the scope: 6,756,355 strict-suffix prefixes from 2,054 certified factor entries reduce the root-4 branch to Phi\_31(269) = sigma(269\textasciicircum\{\}30).
\end{enumerate}

\section{Scope and Open Boundaries}

To prevent over-claiming beyond the proved theorem, we record the explicit boundaries that are \emph{not} decided by the argument above:
\begin{itemize}
  \item nonexistence of odd perfect numbers.
\end{itemize}

\begin{thebibliography}{9}
\bibitem{PrimarySource}
Primary problem list and source literature for Result~8 (elementary and computational number theory), archived with exact theorem locators and verification certificates in the accompanying repository package.
\end{thebibliography}

\end{document}
